\ifdefined\gkver\else\def\gkver{student}\fi
\documentclass[\gkver]{gaokaozhenti}
\begin{document}

\chapter{2014年江苏}

\begin{enumerate}
\item
%% number: 1
%% paperName: 2014年江苏高考第 1 题 3 分
%% typeId: 1
%% type: 单选题
%% chapter: 电磁感应
%% point: 法拉第电磁感应定律
%% method: 无
%% score: 3
%% degree: 650
%% duplicateId: 0
%% body:
如图所示，一正方形线圈的匝数为 $n$，边长为 $a$，线圈平面与匀强磁场垂直，且一半处在磁场中。在 $\Delta t$ 时间内，磁感应强度的方向不变，大小由 $B$ 均匀地增大到 $2B$。在此过程中，线圈中产生的感应电动势为\xzanswer{B}

\onepicture[align=right, w=4.50cm]{figs/01.png}

\fourchoices[answer=B]
{$\frac{Ba^{2}}{2\Delta t}$}
{$\frac{nBa^{2}}{2\Delta t}$}
{$\frac{nBa^{2}}{\Delta t}$}
{$\frac{2nBa^{2}}{\Delta t}$}

%% answer: B
%% memo:
\memoanswer{
本题原卷及材料中未提供详解，待补充。
}

\item
%% number: 2
%% paperName: 2014年江苏高考第 2 题 3 分
%% typeId: 1
%% type: 单选题
%% chapter: 曲线运动
%% point: 人造地球卫星
%% method: 无
%% score: 3
%% degree: 650
%% duplicateId: 0
%% body:
已知地球的质量约为火星质量的 10 倍，地球的半径约为火星半径的 2 倍，则航天器在火星表面附近绕火星做匀速圆周运动的速率约为\xzanswer{A}

\fourchoices[answer=A]
{$3.5\Ukms$}
{$5.0\Ukms$}
{$17.7\Ukms$}
{$35.2\Ukms$}

%% answer: A
%% memo:
\memoanswer{
本题原卷及材料中未提供详解，待补充。
}

\item
%% number: 3
%% paperName: 2014年江苏高考第 3 题 3 分
%% typeId: 1
%% type: 单选题
%% chapter: 交流电
%% point: 变压器
%% method: 无
%% score: 3
%% degree: 650
%% duplicateId: 0
%% body:
远距离输电的原理图如图所示，升压变压器原、副线圈的匝数分别为 $n_{1}$、$n_{2}$，电压分别为 $U_{1}$、$U_{2}$，电流分别为 $I_{1}$、$I_{2}$，输电线上的电阻为 $R$。变压器为理想变压器，则下列关系式中正确的是\xzanswer{D}

\onepicture[align=right, w=7.00cm]{figs/03.png}

\fourchoices[answer=D]
{$\frac{I_{1}}{I_{2}}=\frac{n_{1}}{n_{2}}$}
{$I_{2}=\frac{U_{2}}{R}$}
{$I_{1}U_{1}=I_{2}^{2}R$}
{$I_{1}U_{1}=I_{2}U_{2}$}

%% answer: D
%% memo:
\memoanswer{
本题原卷及材料中未提供详解，待补充。
}

\item
%% number: 4
%% paperName: 2014年江苏高考第 4 题 3 分
%% typeId: 1
%% type: 单选题
%% chapter: 电场
%% point: 电场叠加
%% method: 对称法
%% score: 3
%% degree: 450
%% duplicateId: 0
%% body:
如图所示，一圆环上均匀分布着正电荷，$x$ 轴垂直于环面且过圆心 O。下列关于 $x$ 轴上的电场强度和电势的说法中正确的是\xzanswer{B}

\onepicture[align=right, w=5.50cm]{figs/04.png}

\fourchoices[answer=B]
{O 点的电场强度为零，电势最低}
{O 点的电场强度为零，电势最高}
{从 O 点沿 $x$ 轴正方向，电场强度减小，电势升高}
{从 O 点沿 $x$ 轴正方向，电场强度增大，电势降低}

%% answer: B
%% memo:
\memoanswer{
本题原卷及材料中未提供详解，待补充。
}

\item
%% number: 5
%% paperName: 2014年江苏高考第 5 题 3 分
%% typeId: 1
%% type: 单选题
%% chapter: 直线运动
%% point: 运动图象
%% method: 无
%% score: 3
%% degree: 650
%% duplicateId: 0
%% body:
一汽车从静止开始做匀加速直线运动，然后刹车做匀减速直线运动，直到停止。下列速度 $v$ 和位移 $x$ 的关系图像中，能描述该过程的是\xzanswer{A}

\fourchoices[answer=A, ispicture=true, h=2.50cm]
{figs/05a.png}
{figs/05b.png}
{figs/05c.png}
{figs/05d.png}

%% answer: A
%% memo:
\memoanswer{
本题原卷及材料中未提供详解，待补充。
}

\item
%% number: 6
%% paperName: 2014年江苏高考第 6 题 4 分
%% typeId: 4
%% type: 实验题
%% chapter: 曲线运动
%% point: 平抛运动
%% method: 无
%% score: 4
%% degree: 650
%% duplicateId: 0
%% body:
为了验证平抛运动的小球在竖直方向上做自由落体运动，用如图所示的装置进行实验。小锤打击弹性金属片，A 球水平抛出，同时 B 球被松开，自由下落。关于该实验，下列说法中正确的有\xzanswer{BC}

\onepicture[align=right, h=4.60cm]{figs/06.png}

\fourchoices[answer=BC]
{两球的质量应相等}
{两球应同时落地}
{应改变装置的高度，多次实验}
{实验也能说明 A 球在水平方向上做匀速直线运动}

%% answer: BC
\jdanswer{BC}

%% memo:
\memoanswer{
本题原卷及材料中未提供详解，待补充。
}

\item
%% number: 7
%% paperName: 2014年江苏高考第 7 题 4 分
%% typeId: 2
%% type: 多选题
%% chapter: 交流电
%% point: 交流电
%% method: 无
%% score: 4
%% degree: 450
%% duplicateId: 0
%% body:
如图所示，在线圈上端放置一盛有冷水的金属杯，现接通交流电源，过了几分钟，杯内的水沸腾起来。若要缩短上述加热时间，下列措施可行的有\xzanswer{AB}

\onepicture[align=right, w=4.50cm]{figs/07.png}

\fourchoices[answer=AB]
{增加线圈的匝数}
{提高交流电源的频率}
{将金属杯换为瓷杯}
{取走线圈中的铁芯}

%% answer: AB
%% memo:
\memoanswer{
本题原卷及材料中未提供详解，待补充。
}

\item
%% number: 8
%% paperName: 2014年江苏高考第 8 题 4 分
%% typeId: 2
%% type: 多选题
%% chapter: 运动和力
%% point: 牛顿第二定律
%% method: 整体与隔离
%% score: 4
%% degree: 450
%% duplicateId: 0
%% body:
如图所示，A、B 两物块的质量分别为 $2m$ 和 $m$，静止叠放在水平地面上。A、B 间的动摩擦因数为 $\mu$，B 与地面间的动摩擦因数为 $\frac{1}{2}\mu$。最大静摩擦力等于滑动摩擦力，重力加速度为 $g$。现对 A 施加一水平拉力 $F$，则\xzanswer{BCD}

\onepicture[align=right, w=6.50cm]{figs/08.png}

\fourchoices[answer=BCD]
{当 $F<2\mu mg$ 时，A、B 都相对地面静止}
{当 $F=\frac{5}{2}\mu mg$ 时，A 的加速度为 $\frac{1}{3}\mu g$}
{当 $F>3\mu mg$ 时，A 相对 B 滑动}
{无论 $F$ 为何值，B 的加速度不会超过 $\frac{1}{2}\mu g$}

%% answer: BCD
%% memo:
\memoanswer{
本题原卷及材料中未提供详解，待补充。
}

\item
%% number: 9
%% paperName: 2014年江苏高考第 9 题 4 分
%% typeId: 2
%% type: 多选题
%% chapter: 磁场
%% point: 洛伦兹力
%% method: 无
%% score: 4
%% degree: 450
%% duplicateId: 0
%% body:
如图所示，导电物质为电子的霍尔元件位于两串联线圈之间，线圈中电流为 $I$，线圈间产生匀强磁场，磁感应强度大小 $B$ 与 $I$ 成正比，方向垂直于霍尔元件的两侧面，此时通过霍尔元件的电流为 $I_{H}$，与其前后表面相连的电压表测出的霍尔电压 $U_{H}$ 满足：$U_{H}=k\frac{I_{H}B}{d}$，式中 $k$ 为霍尔系数，$d$ 为霍尔元件两侧面间的距离。电阻 $R$ 远大于 $R_{L}$，霍尔元件的电阻可以忽略，则\xzanswer{CD}

\onepicture[align=right, w=8.00cm]{figs/09.png}

\fourchoices[answer=CD]
{霍尔元件前表面的电势低于后表面}
{若电源的正负极对调，电压表将反偏}
{$I_{H}$ 与 $I$ 成正比}
{电压表的示数与 $R_{L}$ 消耗的电功率成正比}

%% answer: CD
%% memo:
\memoanswer{
本题原卷及材料中未提供详解，待补充。
}

\item
%% number: 10
%% paperName: 2014年江苏高考第 10 题 8 分
%% typeId: 4
%% type: 实验题
%% chapter: 电路
%% point: 电路实验
%% method: 无
%% score: 8
%% degree: 450
%% duplicateId: 0
%% body:
某同学通过实验测量一种合金的电阻率。

\begin{enumerate}
\item
用螺旋测微器测量合金丝的直径。为防止读数时测微螺杆发生转动，读数前应先旋紧图1所示的部件\_\_\_\_\_\_\_\_（选填“A”、“B”、“C”或“D”）。从图中的示数可读出合金丝的直径为\_\_\_\_\_\_\_\_$\Umm$。
\item
如图2所示是测量合金丝电阻的电路，相关器材的规格已在图中标出。合上开关，将滑动变阻器的滑片移到最左端的过程中，发现电压表和电流表的指针只在图示位置发生很小的变化。由此可以推断：电路中\_\_\_\_\_\_（选填图中表示接线柱的数字）之间出现了\_\_\_\_\_\_（选填“短路”或“断路”）。
\item
在电路故障被排除后，调节滑动变阻器，读出电压表和电流表的示数分别为 $2.23\UV$ 和 $38\UmA$，由此，该同学算出接入电路部分的合金丝的阻值为 $58.7\UO$。为了更准确地测出合金丝的阻值，在不更换实验器材的条件下，对实验应作怎样的改进？请写出两条建议。
\end{enumerate}

\twopicture[numA=1, numB=2, labelprefix={图}, align=right, h=3.80cm]
{figs/10a.png}
{figs/10b.png}

%% answer:
\jdanswer{
\begin{enumerate}
\item
B，0.410
\item
7、9；断路
\item
电流表改为内接；测量多组电流和电压值，计算出电阻的平均值。（或测量多组电流和电压值，用图像法求电阻值）
\end{enumerate}
}

%% memo:
\memoanswer{
本题原卷及材料中未提供详解，待补充。
}

\item
%% number: 11
%% paperName: 2014年江苏高考第 11 题 10 分
%% typeId: 4
%% type: 实验题
%% chapter: 力物体的平衡
%% point: 验证平行四边形实验
%% method: 无
%% score: 10
%% degree: 450
%% duplicateId: 0
%% body:
小明通过实验验证力的平行四边形定则。

\begin{enumerate}
\item
实验记录纸如图1所示，O 点为橡皮筋被拉伸后伸长到的位置，两弹簧测力计共同作用时，拉力 $F_{1}$ 和 $F_{2}$ 的方向分别过 $P_{1}$ 和 $P_{2}$ 点，一个弹簧测力计拉橡皮筋时，拉力 $F_{3}$ 的方向过 $P_{3}$ 点。三个力的大小分别为：$F_{1}=3.30\UN$、$F_{2}=3.85\UN$ 和 $F_{3}=4.25\UN$。请根据图中给出的标度作图求出 $F_{1}$ 和 $F_{2}$ 的合力。
\item
仔细分析实验，小明怀疑实验中的橡皮筋被多次拉伸后弹性发生了变化，影响实验结果。他用弹簧测力计先后两次将橡皮筋拉伸到相同长度，发现读数不相同，于是进一步探究了拉伸过程对橡皮筋弹性的影响。

实验装置如图2所示，将一张白纸固定在竖直放置的木板上，橡皮筋的上端固定于 O 点，下端 N 挂一重物。用与白纸平行的水平力缓慢地移动 N，在白纸上记录下 N 的轨迹。重复上述过程，再次记录下 N 的轨迹。

两次实验记录的轨迹如图3所示。过 O 点作一条直线与轨迹交于 $a$、$b$ 两点，则实验中橡皮筋分别被拉伸到 a 和 b 时所受拉力 $F_{a}$、$F_{b}$ 的大小关系为\_\_\_\_\_\_\_\_。
\item
根据（2）中的实验，可以得出的实验结果有哪些？（填写选项前的字母）

（A）橡皮筋的长度与受到的拉力成正比

（B）两次受到的拉力相同时，橡皮筋第 2 次的长度较长

（C）两次被拉伸到相同长度时，橡皮筋第 2 次受到的拉力较大

（D）两次受到的拉力相同时，拉力越大，橡皮筋两次的长度之差越大
\item
根据小明的上述实验探究，请对验证力的平行四边形定则实验提出两点注意事项。
\end{enumerate}

\threepicture[numA=1, numB=2, numC=3, labelprefix={图}, align=right, w=4.20cm]
{figs/11a.png}
{figs/11b.png}
{figs/11c.png}

%% answer:
\jdanswer{
\begin{enumerate}
\item
见解析图，$F_{\text{合}}=4.6\sim4.9\UN$ 都算对。
\item
$F_{a}=F_{b}$。
\item
B、D。
\item
橡皮筋拉伸不宜过长；选用新橡皮筋。（或：拉力不宜过大；选用弹性较好的橡皮筋；换用弹性好的弹簧。）
\end{enumerate}
}

%% memo:
\memoanswer{
（1）用直尺作出如图所示的力的示意图，得到 $F_{\text{合}}=4.6\sim4.9\UN$ 都算对。

\onepicture[w=4.80cm]{figs/11_an.png}

（2）根据力的平衡 $F=mg\tan\theta$，所以 $F_{a}=F_{b}$，$\theta$ 为橡皮筋被拉伸的方向与竖直方向的夹角。

（3）选 B、D。根据题图，两次受到的拉力相同时，橡皮筋第 2 次的长度较长，B 项正确；根据题图，两次被拉伸到相同长度时，即 $Oa=Ob$，根据 $F=mg\tan\theta$，橡皮筋第 2 次受到的拉力较小，C 项错误；从题图可以看出，两次受到的拉力相同时，拉力越大，橡皮筋两次的长度之差越大，D 正确，A 错误。

（4）橡皮筋拉伸不宜过长；选用新橡皮筋。（或拉力不宜过大，选用弹性好的橡皮筋；换用弹性好的弹簧）
}

\item
%% number: 12
%% paperName: 2014年江苏高考第 12 题 5 分
%% typeId: 6
%% type: 计算题
%% chapter: 运动和力
%% point: 动量守恒定律
%% method: 无
%% score: 5
%% degree: 650
%% duplicateId: 0
%% body:
【选做题】本题包括 A、B、C 三小题，请选定其中两小题，并在相应的答题区域内作答。若多做，则按 A、B 两小题评分。

\textbf{A．}【选修 3-3】（12 分）

一种海浪发电机的气室如图所示。工作时，活塞随海浪上升或下降，改变气室中空气的压强，从而驱动进气阀门和出气阀门打开或关闭。气室先后经历吸入、压缩和排出空气的过程，推动出气口处的装置发电。气室中的空气可视为理想气体。

\begin{enumerate}
\item
下列对理想气体的理解，正确的有\xzanswer{AD}

\fourchoices[answer=AD]
{理想气体实际上并不存在，只是一种理想模型}
{只要气体压强不是很高就可视为理想气体}
{一定质量的某种理想气体的内能与温度、体积都有关}
{在任何温度、任何压强下，理想气体都遵循气体实验定律}
\item
压缩过程中，两个阀门均关闭。若此过程中，气室中的气体与外界无热量交换，内能增加了 $3.4\times10^{4}\UJ$，则该气体的分子平均动能\tkanswer{增大}（选填“增大”、“减小”或“不变”），活塞对该气体所做的功\tkanswer{等于}（选填“大于”、“小于”或“等于”）$3.4\times10^{4}\UJ$。
\item
上述过程中，气体刚被压缩时的温度为 $27\UdC$，体积为 $0.224\Umc$，压强为 1 个标准大气压。已知 $1\Umol$ 气体在 1 个标准大气压、$0\UdC$ 时的体积为 $22.4\UL$，阿伏加德罗常数 $N_{A}=6.02\times10^{23}\Umoln$。计算此时气室中气体的分子数。（计算结果保留一位有效数字）
\end{enumerate}

\onepicture[align=right, h=4.20cm]{figs/12a.png}

\textbf{B．}【选修 3-4】（12 分）

\begin{enumerate}
\item
某同学用单色光进行双缝干涉实验，在屏上观察到图（甲）所示的条纹，仅改变一个实验条件后，观察到的条纹如题图（乙）所示。他改变的实验条件可能是\xzanswer{B}

\fourchoices[answer=B]
{减小光源到单缝的距离}
{减小双缝之间的距离}
{减小双缝到光屏之间的距离}
{换用频率更高的单色光源}

\twopicture[numA=甲, numB=乙, labelprefix={图}, w=6.20cm]
{figs/12b.png}
{figs/12c.png}
\item
在“探究单摆的周期与摆长的关系”实验中，某同学准备好相关实验器材后，把单摆从平衡位置拉开一个很小的角度后释放，同时按下秒表开始计时，当单摆再次回到释放位置时停止计时，将记录的这段时间作为单摆的周期。以上操作中有不妥之处，请对其中两处加以改正。
\item
Morpho 蝴蝶的翅膀在阳光的照射下呈现出闪亮耀眼的蓝色光芒，这是因为光照射到翅膀的鳞片上发生了干涉。电子显微镜下鳞片结构的示意图如图所示。一束光以入射角 $i$ 从 a 点入射，经过折射和反射后从 b 点出射。设鳞片的折射率为 $n$，厚度为 $d$，两片之间空气层厚度为 $h$。取光在空气中的速度为 $c$，求光从 a 到 b 所需的时间 $t$。
\end{enumerate}

\onepicture[align=right, w=7.00cm]{figs/12d.png}

\textbf{C．}【选修 3-5】（12 分）

\begin{enumerate}
\item
已知钙和钾的截止频率分别为 $7.73\times10^{14}\UHz$ 和 $5.44\times10^{14}\UHz$，在某种单色光的照射下两种金属均发生光电效应，比较它们表面逸出的具有最大初动能的光电子，钙逸出的光电子具有较大的\xzanswer{A}

\fourchoices[answer=A]
{波长}
{频率}
{能量}
{动量}
\item
氡 222 是一种天然放射性气体，被吸入后，会对人的呼吸系统造成辐射损伤。它是世界卫生组织公布的主要环境致癌物质之一。其衰变方程是 \ce{^{222}_{86}Rn -> ^{218}_{84}Po} $+$ \tkanswer{\ce{^{4}_{2}He}}。已知 \ce{^{222}_{86}Rn} 的半衰期约为 $3.8\Ud$，则约经过\tkanswer{$15.2\Ud$}，$16\Ug$ 的 \ce{^{222}_{86}Rn} 衰变后还剩 $1\Ug$。
\item
牛顿的《自然哲学的数学原理》中记载，A、B 两个玻璃球相碰，碰撞后的分离速度和它们碰撞前的接近速度之比总是约为 $15:16$。分离速度是指碰撞后 B 对 A 的速度，接近速度是指碰撞前 A 对 B 的速度。若上述过程是质量为 $2m$ 的玻璃球 A 以速度 $v_{0}$ 碰撞质量为 $m$ 的静止玻璃球 B，且为对心碰撞，求碰撞后 A、B 的速度大小。
\end{enumerate}

%% answer:
\jdanswer{
\begin{enumerate}
\item
A：（1）AD；（2）增大，等于；（3）$N=5\times10^{24}$（或 $N=6\times10^{24}$）
\item
B：（1）B；（2）①应在摆球通过平衡位置时开始计时；②应测量单摆多次全振动的时间，再计算出周期的测量值；（3）$t=\frac{2n^{2}d}{c\sqrt{n^{2}-\sin^{2}i}}+\frac{2h}{c\cos i}$
\item
C：（1）A；（2）\ce{^{4}_{2}He}（或 $\alpha$），$15.2\Ud$；（3）$v_{1}=\frac{17}{48}v_{0}$，$v_{2}=\frac{31}{24}v_{0}$
\end{enumerate}
}

%% memo:
\memoanswer{
\textbf{A．}（3）设气体在标准状态时的体积为 $V_{1}$，等压过程 $\frac{V}{T}=\frac{V_{1}}{T_{1}}$，气体物质的量 $n=\frac{V_{1}}{V_{0}}$，且分子数 $N=nN_{A}$，解得 $N=\frac{VT_{1}}{V_{0}T}N_{A}$，带入数据得 $N=5\times10^{24}$（或 $N=6\times10^{24}$）。

\textbf{B．}（3）设光在鳞片中的折射角为 $\gamma$，折射定律 $\sin i=n\sin\gamma$，在鳞片中传播路程 $l_{1}=\frac{2d}{\cos\gamma}$，传播速度 $v=\frac{c}{n}$，传播时间 $t_{1}=\frac{l_{1}}{v}$，解得 $t_{1}=\frac{2n^{2}d}{c\sqrt{n^{2}-\sin^{2}i}}$，同理，在空气中传播时间 $t_{2}=\frac{2h}{c\cos i}$，则 $t=t_{1}+t_{2}=\frac{2n^{2}d}{c\sqrt{n^{2}-\sin^{2}i}}+\frac{2h}{c\cos i}$。

\textbf{C．}（3）设 A、B 球碰撞后速度分别为 $v_{1}$ 和 $v_{2}$，由动量守恒定律 $2mv_{0}=2mv_{1}+mv_{2}$，且由题意知 $\frac{v_{2}-v_{1}}{v_{0}}=\frac{15}{16}$，解得 $v_{1}=\frac{17}{48}v_{0}$，$v_{2}=\frac{31}{24}v_{0}$。
}

\item
%% number: 13
%% paperName: 2014年江苏高考第 13 题 15 分
%% typeId: 6
%% type: 计算题
%% chapter: 电磁感应
%% point: 电磁感应能量分析
%% method: 无
%% score: 15
%% degree: 450
%% duplicateId: 0
%% body:
如图所示，在匀强磁场中有一倾斜的平行金属导轨，导轨间距为 $L$，长为 $3d$，导轨平面与水平面的夹角为 $\theta$，在导轨的中部刷有一段长为 $d$ 的薄绝缘涂层。匀强磁场的磁感应强度大小为 $B$，方向与导轨平面垂直。质量为 $m$ 的导体棒从导轨的顶端由静止释放，在滑上涂层之前已经做匀速运动，并一直匀速滑到导轨底端。导体棒始终与导轨垂直，且仅与涂层间有摩擦，接在两导轨间的电阻为 $R$，其他部分的电阻均不计，重力加速度为 $g$。求：

\begin{enumerate}
\item
导体棒与涂层间的动摩擦因数 $\mu$；
\item
导体棒匀速运动的速度大小 $v$；
\item
整个运动过程中，电阻产生的焦耳热 $Q$。
\end{enumerate}

\onepicture[align=right, w=8.00cm]{figs/13.png}

%% answer:
\jdanswer{
\begin{enumerate}
\item
$\mu=\tan\theta$
\item
$v=\frac{mgR\sin\theta}{B^{2}L^{2}}$
\item
$Q=2\mu mgd\cos\theta-\frac{m^{3}g^{2}R^{2}\sin^{2}\theta}{2B^{4}L^{4}}$
\end{enumerate}
}

%% memo:
\memoanswer{
（1）在绝缘涂层上，受力平衡 $mg\sin\theta=\mu mg\cos\theta$，得 $\mu=\tan\theta$。

（2）在光滑导轨上，感应电动势 $E=BLv$，感应电流 $I=\frac{E}{R}$，安培力 $F_{\mathrm{A}}=BIL$，受力平衡 $mg\sin\theta=F_{\mathrm{A}}$，解得 $v=\frac{mgR\sin\theta}{B^{2}L^{2}}$。

（3）摩擦生热 $Q_{f}=\mu mgd\cos\theta$，能量守恒有 $3\mu mgd\cos\theta=Q+Q_{f}+\frac{1}{2}mv^{2}$，解得焦耳热 $Q=2\mu mgd\cos\theta-\frac{m^{3}g^{2}R^{2}\sin^{2}\theta}{2B^{4}L^{4}}$。
}

\item
%% number: 14
%% paperName: 2014年江苏高考第 14 题 16 分
%% typeId: 6
%% type: 计算题
%% chapter: 磁场
%% point: 带电粒子在磁场中的运动
%% method: 无
%% score: 16
%% degree: 450
%% duplicateId: 0
%% body:
某装置用磁场控制带电粒子的运动，工作原理如图所示。装置的长为 $L$，上下两个相同的矩形区域内存在匀强磁场，磁感应强度大小均为 $B$、方向与纸面垂直且相反，两磁场的间距为 $d$。装置右端有一收集板，M、N、P 为板上的三点，M 位于轴线 $OO'$ 上，N、P 分别位于下方磁场的上、下边界上。在纸面内，质量为 $m$、电荷量为 $-q$ 的粒子以某一速度从装置左端的中点射入，方向与轴线成 $30^{\circ}$ 角，经过上方的磁场区域一次，恰好到达 P 点。改变粒子入射速度的大小，可以控制粒子到达收集板上的位置。不计粒子的重力。

\begin{enumerate}
\item
求磁场区域的宽度 $h$；
\item
欲使粒子到达收集板的位置从 P 点移到 N 点，求粒子入射速度的最小变化量 $\Delta v$；
\item
欲使粒子到达 M 点，求粒子入射速度大小的可能值。
\end{enumerate}

\onepicture[align=right, w=9.00cm]{figs/14.png}

%% answer:
\jdanswer{
\begin{enumerate}
\item
$h=\left(\frac{3}{2}L-d\right)\left(1-\frac{\sqrt{3}}{2}\right)$
\item
$\Delta v=\frac{qB}{m}\left(\frac{L}{6}-\frac{\sqrt{3}}{4}d\right)$
\item
$v_{n}=\frac{qB}{m}\left(\frac{L}{n+1}-\sqrt{3}d\right)$，$\left(1\leqslant n<\frac{\sqrt{3}L}{3d}-1\right.$，$n$ 取整数）
\end{enumerate}
}

%% memo:
\memoanswer{
（1）设粒子在磁场中的轨道半径为 $r$。

根据题意 $L=3r\sin30^{\circ}+3d\cos30^{\circ}$，且 $h=r(1-\cos30^{\circ})$，

解得 $h=\left(\frac{3}{2}L-d\right)\left(1-\frac{\sqrt{3}}{2}\right)$。

（2）设改变入射速度后轨道半径为 $r'$。

$m\frac{v^{2}}{r}=qvB$，$m\frac{v'^{2}}{r'}=qv'B$，

由题意 $3r\sin30^{\circ}=4r'\sin30^{\circ}$，得 $\Delta v=v-v'=\frac{qB}{m}\left(\frac{L}{6}-\frac{\sqrt{3}}{4}d\right)$。

（3）设粒子经过上方磁场 $n$ 次。

由题意 $L=(2n+2)d\cos30^{\circ}+(2n+2)r_{n}\sin30^{\circ}$，

且 $m\frac{v_{n}^{2}}{r_{n}}=qv_{n}B$，解得 $v_{n}=\frac{qB}{m}\left(\frac{L}{n+1}-\sqrt{3}d\right)$，$\left(1\leqslant n<\frac{\sqrt{3}L}{3d}-1\right.$，$n$ 取整数）。
}

\item
%% number: 15
%% paperName: 2014年江苏高考第 15 题 16 分
%% typeId: 6
%% type: 计算题
%% chapter: 直线运动
%% point: 运动的合成
%% method: 无
%% score: 16
%% degree: 250
%% duplicateId: 0
%% body:
如图所示，生产车间有两个相互垂直且等高的水平传送带甲和乙，甲的速度为 $v_{0}$。小工件离开甲前与甲的速度相同，并平稳地传到乙上，工件与乙之间的动摩擦因数为 $\mu$。乙的宽度足够大，重力加速度为 $g$。

\begin{enumerate}
\item
若乙的速度为 $v_{0}$，求工件在乙上侧向（垂直于乙的运动方向）滑过的距离 $s$；
\item
若乙的速度为 $2v_{0}$，求工件在乙上刚停止侧向滑动时的速度大小 $v$；
\item
保持乙的速度 $2v_{0}$ 不变，当工件在乙上刚停止滑动时，下一只工件恰好传到乙上，如此反复。若每个工件的质量均为 $m$，除工件与传送带之间摩擦外，其他能量损耗均不计，求驱动乙的电动机的平均输出功率 $\overline{P}$。
\end{enumerate}

\onepicture[align=right, w=6.50cm]{figs/15.png}

%% answer:
\jdanswer{
\begin{enumerate}
\item
$s=\frac{v_{0}^{2}}{2\mu g}$
\item
$v=2v_{0}$
\item
$\overline{P}=\frac{4\sqrt{5}}{5}\mu mgv_{0}$
\end{enumerate}
}

%% memo:
\memoanswer{
（1）摩擦力与侧向的夹角为 $45^{\circ}$，侧向加速度为 $a_{x}=\mu g\cos45^{\circ}$。

匀变速直线运动 $-2a_{x}s=0-v_{0}^{2}$，解得 $s=\frac{v_{0}^{2}}{2\mu g}$。

（2）设 $t=0$ 时摩擦力与侧向的夹角为 $\theta$，侧向、纵向加速度的大小分别为 $a_{x}$、$a_{y}$，

则 $a_{y}/a_{x}=\tan\theta$。

在很小的时间 $\Delta t$ 内，侧向、纵向的速度增量 $\Delta v_{x}=a_{x}\Delta t$，$\Delta v_{y}=a_{y}\Delta t$，

$\Delta v_{y}/\Delta v_{x}=\tan\theta$。

且由题意知 $\tan\theta=v_{y}/v$，则 $v_{y}'/v_{x}'=(v_{y}-\Delta v_{y})/(v-\Delta v_{x})=\tan\theta$，

所以摩擦力方向不变。

则当 $v_{x}'=0$ 时，$v_{y}'=0$，即 $v=2v_{0}$。

（3）工件在乙上滑动时侧向位移为 $x$，沿乙的方向的位移为 $y$，

由题意知 $a_{x}=\mu g\cos\theta$，$a_{y}=\mu g\sin\theta$。

在侧向上 $-2a_{x}x=0-v_{0}^{2}$，纵向有 $2a_{y}x=4v_{0}^{2}-0$，

工件滑动时间为 $t=2v_{0}/a_{y}$，乙前进的距离为 $y_{1}=2v_{0}t$，

工件相对乙的位移 $L=\sqrt{x^{2}+(y_{1}-y)^{2}}$，则系统产生的热量为 $Q=\mu mgL$，

电动机做功为 $W=\frac{1}{2}m(2v_{0})^{2}-\frac{1}{2}mv_{0}^{2}+Q$，

解得 $\overline{P}=\frac{W}{t}=\frac{4\sqrt{5}}{5}\mu mgv_{0}$。
}

\end{enumerate}

\end{document}
